% =====================================================================================================
% vygotskian_conclusion.tex — The Vygotskian machine, part 6: conclusion
% DRAFT (2026-11-20). Added after the critical review. Include after vygotskian_experiments.tex.
% =====================================================================================================

\section{Conclusion}
\label{sec:conclusion}

We have proposed a way of giving a robot the knowledge contained in a vision-language model without letting that
model write the robot's beliefs. On an existing Active Inference architecture we add an inner voice whose questions
perception answers, a demand signal that decides when the voice is needed, and a cultural plane in which the VLM
supplies room functions, missions and imagination, always as proposals that the robot's experience progressively
replaces.

The machine is inspired by Vygotsky rather than derived from him, and its contribution is the integration of known
components into one probabilistic hierarchy under one rule of evidence. Its Vygotskian claim rests on P7, that a VLM
deprived of the culture attached to its words falls behind in proportion to novelty, and P9 tests the weaker claim that
plan reuse is selected by the robot's uncertainty. Open parts remain. The demand thresholds are declared rather than
derived, the weak grounding constraint lets the VLM's description of an unfamiliar object shape the answer, and a
confidently wrong model stays silent until its errors have consequences, which the experiments will weigh.
